\section{Conclusion \& Future Work} 
One of the practical bottlenecks to long-horizon autonomy is surviving hardware faults without stopping. Building fail-active robots is necessary to maintain functionality under faults without constant human intervention. To that end, we presented \textit{DEFT}, a \textit{D}iffusion-based \textit{E}mbodiment-aware \textit{F}ail-active \textit{T}ask-conditioned trajectory generation framework that enables robots to adapt to joint failures while preserving task performance.
% By treating failure adaptation as conditional trajectory generation \emph{with diffusion}.
Without explicit replanning or policy switching, DEFT generalizes across multiple failure types and motion primitives. Experiments show DEFT maintains high success rates under multiple joint failures, especially in out-of-distribution scenarios. In real-world trials, DEFT executes reliably under previously unseen multi-joint failures, delivering high task success with strict joint limit compliance.
Future work will focus on real-time failure detection, cross-embodiment transfer to enable fail-active strategies learned on one robot architecture to generalize to different robot platforms, and expanding manipulation skills like pivoting or throwing. By reframing how robots respond to actuation failures, DEFT paves the way to the development of reliable autonomous systems capable of continuous operation under adverse conditions. 
\vspace{-10pt}
% \gm{add cross-embodiment}


% While many robotic systems offer remarkable manipulation capabilities, this mechanical complexity inherently multiplies the possibilities for component degradation and failure. This work proposes DEFT, a Diffusion-based Embodiment-aware Fail-active Task-conditioned trajectory generation framework that enables robots to adapt to arbitrary joint failures while maintaining task capabilities. By framing failure adaptation as a conditional trajectory generation problem and using diffusion models, our approach generalizes across diverse failure types and motion primitives without requiring explicit replanning or policy switching. Our experiments demonstrate DEFT's effectiveness over baseline methods, maintaining high success rates even with multiple simultaneous joint failures across both in-distribution and out-of-distribution conditions. Our future work will focus on developing an online adaptation method that can detect and respond to evolving failure conditions in real-time, rather than relying on predetermined failure states. We will also explore cross-embodiment transfer learning to enable fail-active strategies learned on one robot architecture to generalize to different robot platforms. Finally, we plan to incorporate additional manipulation capabilities such as pivoting or throwing which will further enhance the robot's ability to complete diverse tasks. By reframing how robots respond to mechanical failures, DEFT contributes to the development of reliable autonomous systems capable of continuous operation under adverse conditions. 